\subsection*{运动会补充练习}
\begin{enumerate}
    \item (2025$\cdot$嘉定区高二下期末$\cdot$20)如图,直线 $y = {kx} + b\left( {k > 0,b > 0}\right)$ 与椭圆 $\dfrac{{x}^{2}}{4} + {y}^{2} = 1$ 交于 $A\text{ 、 }B$ 两点,记 $\bigtriangleup {AOB}$ 的面积为 $S$.

    \begin{center}
    \includegraphics[width=0.2\textwidth]{bodyxb/src/0_92_183_382_263_0.jpg}
    \end{center}

    \begin{enumerate}
        \item 当 $k = 1$ ， $b = 1$ 时，求 $S$
        \item 当 $\left| {AB}\right|  = 2$ ， $S = 1$ 时，求直线 ${AB}$ 的方程
        \item 求 $S$ 的最大值
    \end{enumerate}
    \begin{jieda}
        \begin{dingautolist}{172}
            \item 通解法
            \begin{enumerate}[(1)]
                \item 设 $A\left( {{x}_{1},{y}_{1}}\right) ,B\left( {{x}_{2},{y}_{2}}\right)$ . 由 $\left\{  \begin{array}{l} \dfrac{{x}^{2}}{4} + {y}^{2} = 1, \\  y = x + 1, \end{array}\right.$ 得 $\dfrac{5{x}^{2}}{4} + {2x} = 0$ ，解得 ${x}_{1} = 0,{x}_{2} =  - \dfrac{8}{5}$ ，则 $\left| {{x}_{1} - {x}_{2}}\right|  = \dfrac{8}{5}$ . 因此 $S = \dfrac{1}{2}\left| b\right| \left| {{x}_{1} - {x}_{2}}\right|  = \dfrac{4}{5}$ .

                \item 由 $\left\{  \begin{array}{l} \dfrac{{x}^{2}}{4} + {y}^{2} = 1, \\  y = {kx} + b, \end{array}\right.$ 得 $\left( {\dfrac{1}{4} + {k}^{2}}\right) {x}^{2} + {2kbx} + {b}^{2} - 1 = 0$ .

                由 $\Delta  = {\left( 2kb\right) }^{2} - 4 \cdot  \left( {\dfrac{1}{4} + {k}^{2}}\right)  \cdot  \left( {{b}^{2} - 1}\right)  = 4{k}^{2} - {b}^{2} + 1 > 0$ , 可得 $\left\{\begin{array}{l} {x}_{1} + {x}_{2} =  - \dfrac{8kb}{1 + 4{k}^{2}} \\  {x}_{1}{x}_{2} = \dfrac{4{b}^{2} - 4}{1 + 4{k}^{2}} \end{array} \right.$

                由 $\left| {AB}\right|  = \left| {{x}_{1} - {x}_{2}}\right| \sqrt{1 + {k}^{2}} = 2,S = \dfrac{\left| b\right| }{2}\left| {{x}_{1} - {x}_{2}}\right|  = 1$ 得 $1 + {k}^{2} = {b}^{2}$ .

                又 $\left| {{x}_{1} - {x}_{2}}\right|  = \sqrt{{\left( {x}_{1} + {x}_{2}\right) }^{2} - 4{x}_{1}{x}_{2}} = \dfrac{4\sqrt{4{k}^{2} + 1 - {b}^{2}}}{1 + 4{k}^{2}}$ ,

                代入 $\left| {AB}\right|  = \left| {{x}_{1} - {x}_{2}}\right| \sqrt{1 + {k}^{2}} = 2$ ，得 $4{k}^{4} - 4{k}^{2} + 1 = 0$ ，解得 $k = \dfrac{\sqrt{2}}{2}$ ， $b = \dfrac{\sqrt{6}}{2}$ .

                直线 ${AB}$ 的方程是 $y = \dfrac{\sqrt{2}}{2}x + \dfrac{\sqrt{6}}{2}$ .

                \item $S = \dfrac{\left| b\right| }{2}\left| {{x}_{1} - {x}_{2}}\right|  = \dfrac{\left| b\right| }{2} \cdot  \dfrac{4\sqrt{4{k}^{2} + 1 - {b}^{2}}}{1 + 4{k}^{2}} = \dfrac{2\left| b\right| \sqrt{4{k}^{2} + 1 - {b}^{2}}}{4{k}^{2} + 1}$ .

                由基本不等式得 $\dfrac{2\left| b\right| \sqrt{4{k}^{2} + 1 - {b}^{2}}}{4{k}^{2} + 1} \leq  \dfrac{4{k}^{2} + 1 - {b}^{2} + {b}^{2}}{4{k}^{2} + 1} = 1$ ,

                当且仅当 $2{b}^{2} = 4{k}^{2} + 1$ ,等号成立.

                由第(2)小题的结论知, $S$ 可以取到 1,因此 $S$ 的最大值是 1.
            \end{enumerate}
            \item 三角换元
            \begin{enumerate}[(1)]
                \item 同上
                \item 设$A(2cos\alpha, sin\alpha)$，$B(2cos\beta, sin\beta)$，$1 = S = \dfrac{1}{2}|2sin(\alpha - \beta)|$，由于$A, B$的轮换对称性，不妨设$\beta = \alpha + \dfrac{\pi}{2}$，因为$k > 0, b > 0$，故应有$\alpha \in \left(\dfrac{\pi}{4}, \dfrac{3\pi}{4}\right)$
                
                故根据诱导公式有$B(-2sin\alpha, cos\alpha)$，$4 = |AB|^2 = 4(sin\alpha + cos\alpha)^2 + (sin\alpha - cos\alpha)^2 = 5 + 3 sin2\alpha$，因此$sin2\alpha = -\dfrac{1}{3}$，根据$\alpha \in \left(\dfrac{\pi}{4}, \dfrac{3\pi}{4}\right)$可知$cos 2\alpha = -\dfrac{2\sqrt{2}}{3}$，故$tan\alpha = \dfrac{1-cos2\alpha}{sin2\alpha} = -(3+2\sqrt{2})$

                故$k = \dfrac{sin\alpha - cos\alpha}{2(cos\alpha+sin\alpha)} = \dfrac{1}{2} \cdot \dfrac{tan\alpha - 1}{1+tan\alpha} = \dfrac{\sqrt{2}}{2}$，$b = sin\alpha - \sqrt{2}cos\alpha = -3(1+\sqrt{2})cos\alpha = 3(1+\sqrt{2}) \times \sqrt{\dfrac{1}{1+(3+2\sqrt{2})^2}} = 3\sqrt{\dfrac{3+2
                \sqrt{2}}{18+12\sqrt{2}}} = \dfrac{\sqrt{6}}{2}$

                直线 ${AB}$ 的方程是 $y = \dfrac{\sqrt{2}}{2}x + \dfrac{\sqrt{6}}{2}$.
                \item $S = |sin(\alpha - \beta)| \leq 1$， 由第(2)小题的结论知, $S$ 可以取到 1,因此 $S$ 的最大值是 1.
            \end{enumerate}
        \end{dingautolist}
    \end{jieda}
    \item (2008$\cdot$安徽卷（理）$\cdot$18)如图,在四棱锥 $O - {ABCD}$ 中,底面 ${ABCD}$ 四边长为 1 的菱形, $\angle {ABC} = \dfrac{\pi }{4},{OA} \bot$ 底面 ${ABCD},{OA} = 2,M$ 为 ${OA}$ 的中点， $N$ 为 ${BC}$ 的中点.
    \begin{center}
    \includegraphics[width=0.15\textwidth]{bodyxb/src/14_1349_456_214_246_0.jpg}
    \end{center}
    \begin{enumerate}
        \item 证明:直线 ${MN}//$ 面 ${OCD}$
        \item 求异面直线 ${AB}$ 与 ${MD}$ 所成角的大小
        \item 求点 $B$ 到平面 ${OCD}$ 的距离
    \end{enumerate}
    \begin{jieda}
        \begin{enumerate}
            \item 过 $A$ 作 ${AP} \bot  {CD}$ 交 ${CD}$ 于点 $P$ . 如图示,分别以 $\vv{AB},\vv{AP},A\vv{O}$ 为 $x\text{ 、 }y\text{ 、 }z$ 轴正方向建立坐标系,则 $A\left( {0,0,0}\right) ,B\left( {1,0,0}\right)$ ,

            \begin{center}
            \includegraphics[width=0.2\textwidth]{bodyxb/src/14_1234_810_208_244_0.jpg}
            \end{center}


            $P\left( {0,\dfrac{\sqrt{2}}{2},0}\right) ,D\left( {-\dfrac{\sqrt{2}}{2},\dfrac{\sqrt{2}}{2},0}\right) ,O\left( {0,0,2}\right) ,M\left( {0,0,1}\right) ,N\left( {1 - \dfrac{\sqrt{2}}{4},\dfrac{\sqrt{2}}{4},0}\right) .$

            $\vv{MN} = \left( {1 - \dfrac{\sqrt{2}}{4},\dfrac{\sqrt{2}}{4}, - 1}\right) ,\vv{OP} = \left( {0,\dfrac{\sqrt{2}}{2}, - 2}\right) ,\vv{OD} = \left( {-\dfrac{\sqrt{2}}{2},\dfrac{\sqrt{2}}{2}, - 2}\right)$ .

            设平面 ${OCD}$ 的法向量为 $\vv{n} = \left( {x,y,z}\right)$ ,则 $\left\{  \begin{array}{l} \vv{n} \cdot  \vv{OP} = \dfrac{\sqrt{2}}{2}y - {2z} = 0 \\  \vv{n} \cdot  \vv{OD} =  - \dfrac{\sqrt{2}}{2}x + \dfrac{\sqrt{2}}{2}y - {2z} = 0 \end{array}\right.$ . 不妨取 $z = \sqrt{2}$ ,解得: $\vv{n} = \left( {0,4,\sqrt{2}}\right)$ . 因为 $\vv{MN} \cdot  \vv{n} = \left( {1 - \dfrac{\sqrt{2}}{4},\dfrac{\sqrt{2}}{4}, - 1}\right)  \cdot  \left( {0,4,\sqrt{2}}\right)  = 0 + \sqrt{2} - \sqrt{2} = 0$ ,直线 ${MN} \not \subset$ 面 ${OCD}$ ,所以 ${MN}//$ 面 ${OCD}$ .

            \item 设直线 ${AB}$ 与 ${MD}$ 所成角为 $\theta$ ，则 $\theta  \in  \left( {0,\dfrac{\pi }{2}}\right\rbrack$ . 因为 $\vv{AB} = \left( {1,0,0}\right)$ ， $\vv{MD} = \left( {-\dfrac{\sqrt{2}}{2},\dfrac{\sqrt{2}}{2}, - 1}\right)$ ， 所以 $\cos \theta  = \left| {\cos \langle \vv{AB},\vv{MD}\rangle }\right|  = \left| \dfrac{\vv{AB} \cdot  \vv{MD}}{\left| \vv{AB}\right|  \times  \left| \vv{MD}\right| }\right|  = \left| \dfrac{-\dfrac{\sqrt{2}}{2} + 0 + 0}{1 \times  \sqrt{\dfrac{1}{2} + \dfrac{1}{2} + 1}}\right|  = \dfrac{1}{2}$ ，所以 $\theta  = \dfrac{\pi }{3}$ ，即直线 ${AB}$ 与 ${MD}$ 所成角为 $\dfrac{\pi }{3}$ .

            \item 设点 $B$ 到平面 ${OCD}$ 的距离为 $d$ ，则 $d$ 为 $\vv{OB}$ 在向量 $\vv{n} = \left( {0,4,\sqrt{2}}\right)$ 上的投影的绝对值， 由 $\vv{OB} = \left( {1,0, - 2}\right)$ ,得 $d = \dfrac{\left| \vv{OB} \cdot  \vv{n}\right| }{\left| \vv{n}\right| } = \dfrac{\left| 0 + 0 - 2\sqrt{2}\right| }{\sqrt{0 + {4}^{2} + 2}} = \dfrac{2}{3}$. 所以点 $B$ 到平面 ${OCD}$ 的距离为 $\dfrac{2}{3}$.
        \end{enumerate}
    \end{jieda}
\end{enumerate}